\chapter{Introducción}

\section{Objetivo}
El objetivo principal del proyecto es desarrollar un sistema que permita analizar el comportamiento de los usuarios mediante el uso de un \emph{score} dinámico que evalúe la desviación respecto a la conducta habitual de cada usuario y, en función de ello, tomar un conjunto de decisiones de forma inmediata.

\section{Alcance}
El proyecto consiste en el desarrollo de un agente autónomo UEBA que trabaja en conjunto con una API. En esta última se implementará un motor de Inteligencia Artificial entrenado por los alumnos para detectar anomalías en el comportamiento de los usuarios. Finalmente, el sistema contará con un \emph{dashboard} que permitirá visualizar el nivel de confianza de cada usuario y modificar los umbrales que definen las acciones a realizar.

\section{Descripción}
Los sistemas de seguridad actuales suelen confiar en el usuario solo en el momento del inicio de sesión. Esto deja un punto ciego crítico: si un atacante logra robar credenciales válidas, puede moverse dentro de la red sin ser detectado. Este problema es cada vez más relevante, ya que las soluciones tradicionales han demostrado quedarse limitadas frente a amenazas más sofisticadas \parencite{Ansari2025}. Ante este escenario, el proyecto propone sumar una capa de monitoreo continuo que acompañe toda la sesión. Esta solución actúa como un complemento a los sistemas de gestión de identidad (IAM) tradicionales y puede integrarse con ellos para validar no solo quién accede, sino también cómo se comporta. La idea es pasar a un modelo capaz de adaptarse a comportamientos cambiantes y detectar anomalías de forma más precisa.  

Para lograrlo, el sistema se organiza en tres componentes: un agente instalado en el equipo local, una API central desplegada en la nube y un panel de administración (\emph{dashboard}). El agente es un programa liviano que se ejecuta en segundo plano y recolecta información sobre la actividad del usuario, como procesos ejecutados e intentos de acceso a sistemas, aplicaciones o recursos de red. Estos datos se envían de forma segura a la API, que calcula en tiempo real el nivel de confianza mediante un modelo de inteligencia artificial y, con base en ese resultado, define la acción a tomar.

Dicho nivel de confianza no se calcula únicamente por la sesión; por el contrario, se apoya en el historial previo del usuario, manteniendo una continuidad entre sesiones, pero aplicando un decaimiento en el factor de riesgo con el paso del tiempo. En función del nivel de riesgo detectado, la API devuelve una respuesta al agente indicando la acción a ejecutar, que puede variar desde una simple alerta hasta restricciones más severas como limitación de permisos.

Por último, el sistema incluye un \emph{dashboard web} que permite visualizar el estado de los usuarios y sus niveles de riesgo en tiempo real. Desde allí, los administradores pueden ajustar los umbrales de confianza y definir qué acciones se ejecutan en cada caso, adaptando el sistema a distintos contextos de seguridad. Además, el modelo de inteligencia artificial se reentrena para incorporar nuevos comportamientos de los usuarios, lo que permite mantener su precisión a lo largo del tiempo. El periodo en el que se realiza este proceso puede configurarse desde el panel de administración, de modo que no interfiera con el funcionamiento del entorno productivo.

% Puedes descomentar y adaptar este párrafo según los requerimientos de tus docentes:
% El presente documento se estructura de la siguiente manera: en el Capítulo 2 se detalla el Marco Teórico, en el Capítulo 3 se presenta el Diseño de la Solución...